\documentclass[12pt,a4paper]{report}
\usepackage[utf8]{inputenc}
\usepackage[T1]{fontenc}
\usepackage[french]{babel}
\usepackage{geometry}
\geometry{top=2.5cm,bottom=2.5cm,left=2.5cm,right=2.5cm}
\usepackage{graphicx}
\usepackage{xcolor}
\usepackage{tikz}
\usetikzlibrary{shapes.geometric,arrows.meta,positioning,calc,shadows.blur,decorations.pathmorphing,fit,backgrounds,patterns}
\usepackage{tcolorbox}
\tcbuselibrary{skins,breakable}
\usepackage{fontawesome5}
\usepackage{enumitem}
\usepackage{hyperref}
\hypersetup{colorlinks=true,linkcolor=indigo,urlcolor=indigo,pdfauthor={MonEcole},pdftitle={Manuel — Module Évaluations}}
\usepackage{fancyhdr}
\usepackage{tabularx}
\usepackage{booktabs}
\usepackage{multicol}
\usepackage{pifont}
\usepackage{setspace}
\usepackage{wrapfig}
\usepackage{float}
\usepackage{caption}
\usepackage{mdframed}
\usepackage{lipsum}

% ====== COULEURS ======
\definecolor{indigo}{HTML}{4338CA}
\definecolor{violet}{HTML}{7C3AED}
\definecolor{amber}{HTML}{D97706}
\definecolor{emerald}{HTML}{059669}
\definecolor{rose}{HTML}{E11D48}
\definecolor{slate}{HTML}{334155}
\definecolor{slatelight}{HTML}{94A3B8}
\definecolor{teal}{HTML}{0D9488}
\definecolor{sky}{HTML}{0EA5E9}
\definecolor{coverbg}{HTML}{1E1B4B}
\definecolor{coveraccent}{HTML}{6366F1}
\definecolor{lightbg}{HTML}{F8FAFC}
\definecolor{warningbg}{HTML}{FEF3C7}
\definecolor{successbg}{HTML}{ECFDF5}
\definecolor{infobg}{HTML}{EFF6FF}
\definecolor{dangerbg}{HTML}{FEF2F2}
\definecolor{orangebg}{HTML}{FFF7ED}
\definecolor{purplebg}{HTML}{FAF5FF}
\definecolor{tealbg}{HTML}{F0FDFA}

% ====== TCOLORBOX STYLES ======
\newtcolorbox{astuce}[1][]{enhanced,colback=successbg,colframe=emerald,fonttitle=\bfseries\small,title={\faLightbulb\hspace{6pt}Astuce},left=8pt,right=8pt,top=6pt,bottom=6pt,arc=4pt,boxrule=1pt,breakable,shadow={2pt}{-2pt}{0pt}{emerald!20},#1}
\newtcolorbox{attention}[1][]{enhanced,colback=warningbg,colframe=amber,fonttitle=\bfseries\small,title={\faExclamationTriangle\hspace{6pt}Attention},left=8pt,right=8pt,top=6pt,bottom=6pt,arc=4pt,boxrule=1pt,breakable,shadow={2pt}{-2pt}{0pt}{amber!20},#1}
\newtcolorbox{info}[1][]{enhanced,colback=infobg,colframe=indigo,fonttitle=\bfseries\small,title={\faInfoCircle\hspace{6pt}Information},left=8pt,right=8pt,top=6pt,bottom=6pt,arc=4pt,boxrule=1pt,breakable,shadow={2pt}{-2pt}{0pt}{indigo!15},#1}
\newtcolorbox{danger}[1][]{enhanced,colback=dangerbg,colframe=rose,fonttitle=\bfseries\small,title={\faTimesCircle\hspace{6pt}Important},left=8pt,right=8pt,top=6pt,bottom=6pt,arc=4pt,boxrule=1pt,breakable,shadow={2pt}{-2pt}{0pt}{rose!15},#1}
\newtcolorbox{etape}[1][]{enhanced,colback=white,colframe=indigo!60,fonttitle=\bfseries\small,left=10pt,right=10pt,top=8pt,bottom=8pt,arc=6pt,boxrule=1.5pt,breakable,shadow={2pt}{-2pt}{0pt}{black!12},#1}
\newtcolorbox{prerequis}[1][]{enhanced,colback=purplebg,colframe=violet,fonttitle=\bfseries\small,title={\faClipboard\hspace{6pt}Prérequis},left=8pt,right=8pt,top=6pt,bottom=6pt,arc=4pt,boxrule=1pt,breakable,#1}
\newtcolorbox{resultat}[1][]{enhanced,colback=tealbg,colframe=teal,fonttitle=\bfseries\small,title={\faCheck\hspace{6pt}Résultat attendu},left=8pt,right=8pt,top=6pt,bottom=6pt,arc=4pt,boxrule=1pt,breakable,#1}

% ====== HEADERS / FOOTERS ======
\pagestyle{fancy}
\fancyhf{}
\fancyhead[L]{\small\color{slatelight}\leftmark}
\fancyhead[R]{\small\color{indigo}\textbf{MonEcole}}
\fancyfoot[C]{\small\color{slatelight}\thepage}
\renewcommand{\headrulewidth}{0.5pt}
\renewcommand{\footrulewidth}{0.3pt}

% ====== COMMANDES ======
\newcommand{\menubtn}[1]{\tikz[baseline=(X.base)]{\node[fill=indigo!12,rounded corners=3pt,inner sep=3pt,font=\small\sffamily\bfseries\color{indigo}](X){#1};}}
\newcommand{\champ}[1]{\texttt{\color{violet}#1}}
\newcommand{\bouton}[1]{{\colorbox{emerald}{\small\sffamily\bfseries\color{white}\,#1\,}}}
\newcommand{\onglet}[1]{{\colorbox{indigo}{\small\sffamily\bfseries\color{white}\,#1\,}}}
\newcommand{\btnviolet}[1]{{\colorbox{violet}{\small\sffamily\bfseries\color{white}\,#1\,}}}
\newcommand{\btnamber}[1]{{\colorbox{amber}{\small\sffamily\bfseries\color{white}\,#1\,}}}
\newcommand{\btnrose}[1]{{\colorbox{rose}{\small\sffamily\bfseries\color{white}\,#1\,}}}

\begin{document}

% ===================================================================
%                         PAGE DE COUVERTURE
% ===================================================================
\begin{titlepage}
\begin{tikzpicture}[remember picture,overlay]
  \fill[coverbg] (current page.south west) rectangle (current page.north east);
  \fill[coveraccent,opacity=0.12] (current page.north east) ++(-4,1) circle (8cm);
  \fill[violet,opacity=0.08] (current page.south west) ++(3,-2) circle (10cm);
  \fill[coveraccent,opacity=0.06] (current page.center) ++(6,-4) circle (5cm);
  \fill[amber,opacity=0.05] (current page.north west) ++(2,-8) circle (4cm);
  \draw[coveraccent,line width=2pt,opacity=0.4] ([yshift=-6cm]current page.north west) ++(2.5,0) -- ++(16,0);
  \node[fill=coveraccent,rounded corners=14pt,minimum size=2cm,inner sep=0pt] at ([yshift=-3.5cm]current page.north) {\color{white}\fontsize{28pt}{28pt}\selectfont\faFile};
  \node[anchor=north,text width=14cm,align=center] at ([yshift=-6.5cm]current page.north) {
    \color{white}\fontsize{32pt}{36pt}\selectfont\textbf{Manuel d'Utilisation}\\[12pt]
    \fontsize{18pt}{22pt}\selectfont\color{coveraccent!70!white}Module \textbf{Évaluations}
  };
  \node[anchor=north,text width=14cm,align=center] at ([yshift=-10cm]current page.north) {
    \color{white!60}\fontsize{11pt}{15pt}\selectfont
    Guide complet et détaillé pour la gestion des types d'évaluations,\\
    la planification, la saisie des notes, le calcul des moyennes pondérées,\\
    les délibérations et la génération des bulletins scolaires.
  };
  \draw[coveraccent,line width=1.5pt] ([yshift=-12.5cm]current page.north) ++(-3,0) -- ++(6,0);
  \node[anchor=north,text width=14cm,align=center] at ([yshift=-13.5cm]current page.north) {
    \color{white!50}\small
    \faBuilding\hspace{4pt}Plateforme MonEcole — Architecture Hub-and-Spoke\\[6pt]
    \faCalendar\hspace{4pt}Avril 2026 — Version 1.0\\[6pt]
    \faBook\hspace{4pt}Documentation Utilisateur — 4 Onglets Détaillés
  };
  \node[anchor=south,text width=14cm,align=center] at ([yshift=1.5cm]current page.south) {
    \color{white!30}\footnotesize \faCopyright\hspace{3pt} Nexora Tech — Tous droits réservés
  };
\end{tikzpicture}
\end{titlepage}

\tableofcontents
\newpage

% ===================================================================
\chapter{Introduction au Module Évaluations}
% ===================================================================

\section{Présentation générale}

Le module \textbf{Évaluations} est le cœur académique de la plateforme \textbf{MonEcole}. Il couvre l'intégralité du cycle de vie des notes scolaires~: de la création des catégories d'évaluations jusqu'à la génération des bulletins PDF en passant par la saisie, la pondération et la délibération.

Ce module est conçu pour être utilisé par les \textbf{directeurs d'établissement}, les \textbf{préfets d'études} et les \textbf{enseignants}. Il s'adapte automatiquement à la structure de votre pays (cycles, sections, périodes) grâce à l'architecture Hub-and-Spoke de MonEcole.

\begin{info}
Ce module est accessible depuis la barre de navigation latérale (sidebar) en cliquant sur l'icône \faFile\ suivi du texte \textbf{Évaluations}. La page qui s'ouvre comporte plusieurs sections dont la principale est la section \textbf{Évaluations} avec ses 4 onglets.
\end{info}

\section{Les quatre onglets}

La section Évaluations est structurée en \textbf{4 onglets} (pills) situés en haut de la page. Chaque onglet correspond à une étape logique du processus d'évaluation~:

\begin{center}
\begin{tikzpicture}[
  pill/.style={rounded corners=10pt, minimum height=1.1cm, minimum width=4cm, font=\small\bfseries, inner sep=8pt, text=white, drop shadow={shadow xshift=1pt, shadow yshift=-1pt, opacity=0.2}},
  desc/.style={text width=3.8cm, align=center, font=\tiny\color{slate}},
  num/.style={circle, fill=white, text=#1, font=\tiny\bfseries, inner sep=2pt, minimum size=14pt}
]
  \node[pill, fill=indigo] (t1) at (0,0) {\faTags\ Types d'Évaluations};
  \node[pill, fill=teal, right=6pt of t1] (t2) {\faCalendar\ Planifier};
  \node[pill, fill=violet, right=6pt of t2] (t3) {\faPen\ Notes Par Éval.};
  \node[pill, fill=amber, right=6pt of t3] (t4) {\faBalanceScale\ Notes Pondérées};

  \node[num=indigo, above left=2pt and -8pt of t1] {1};
  \node[num=teal, above left=2pt and -8pt of t2] {2};
  \node[num=violet, above left=2pt and -8pt of t3] {3};
  \node[num=amber, above left=2pt and -8pt of t4] {4};

  \node[desc, below=8pt of t1] {Définir les catégories d'évaluations (IE, DS, TP, EX...)};
  \node[desc, below=8pt of t2] {Créer et planifier chaque évaluation avec date et barème};
  \node[desc, below=8pt of t3] {Saisir les notes brutes de chaque élève};
  \node[desc, below=8pt of t4] {Calculer TJ, EX et Totaux pondérés};
\end{tikzpicture}
\end{center}

\section{Flux de travail complet}

Le processus d'évaluation suit un flux séquentiel strict. Chaque étape dépend de la précédente~:

\begin{center}
\begin{tikzpicture}[
  step/.style={draw=#1!60, fill=#1!8, rounded corners=8pt, minimum width=2.4cm, minimum height=1.5cm, align=center, font=\scriptsize\bfseries, text=slate, drop shadow={shadow xshift=1pt, shadow yshift=-1pt, opacity=0.12}},
  arrow/.style={-{Stealth[length=6pt,width=5pt]}, thick, color=#1!60},
  lbl/.style={font=\tiny\itshape\color{slatelight}, above, midway}
]
  \node[step=indigo] (s1) at (0,0) {\faTags\\[2pt]Créer les\\Types};
  \node[step=teal, right=0.6cm of s1] (s2) {\faCalendar\\[2pt]Planifier les\\Évaluations};
  \node[step=violet, right=0.6cm of s2] (s3) {\faPen\\[2pt]Saisir les\\Notes brutes};
  \node[step=amber, right=0.6cm of s3] (s4) {\faCalculator\\[2pt]Calculer\\TJ / EX};
  \node[step=emerald, right=0.6cm of s4] (s5) {\faGavel\\[2pt]Délibérer};
  \node[step=rose, right=0.6cm of s5] (s6) {\faFile\\[2pt]Générer\\Bulletins};

  \draw[arrow=indigo] (s1) -- (s2) node[lbl]{types};
  \draw[arrow=teal] (s2) -- (s3) node[lbl]{évals};
  \draw[arrow=violet] (s3) -- (s4) node[lbl]{notes};
  \draw[arrow=amber] (s4) -- (s5) node[lbl]{totaux};
  \draw[arrow=emerald] (s5) -- (s6) node[lbl]{résultats};
\end{tikzpicture}
\end{center}

\begin{danger}
Il est \textbf{impératif} de suivre cet ordre. Par exemple, vous ne pouvez pas saisir de notes (onglet 3) si aucune évaluation n'a été planifiée (onglet 2). De même, les calculs de l'onglet 4 nécessitent que des notes brutes aient été saisies à l'onglet 3.
\end{danger}

\section{Concepts clés à comprendre}

Avant de plonger dans chaque onglet, voici les concepts fondamentaux~:

\begin{description}[leftmargin=1.5cm, style=nextline, font=\bfseries\color{indigo}]
  \item[Type d'évaluation] Catégorie générique (ex.\ : Interrogation Écrite, Devoir Surveillé). Partagé au niveau national.
  \item[Évaluation] Instance concrète d'un type, rattachée à un cours et une période, avec date et barème.
  \item[Répartition] Découpage temporel de l'année (Trimestres $\rightarrow$ Périodes, ou Semestres $\rightarrow$ Périodes).
  \item[Période] Subdivision d'un trimestre/semestre (ex.\ : 1\textsuperscript{ère} Période, 2\textsuperscript{ème} Période).
  \item[Note Pondérée (NP)] Moyenne calculée des notes brutes d'une période, pondérée par les poids de chaque évaluation.
  \item[TJ (Travaux Journaliers)] Somme des NP de toutes les périodes d'un semestre/trimestre.
  \item[EX (Examen)] Note d'examen semestriel/trimestriel.
  \item[TOTAL] Somme finale~: TJ + EX pour chaque cours par semestre.
\end{description}

\begin{center}
\begin{tikzpicture}[
  level/.style={draw=#1!50,fill=#1!6,rounded corners=6pt,minimum width=3.5cm,minimum height=0.8cm,font=\scriptsize\bfseries,text=#1,align=center},
  child/.style={draw=#1!40,fill=#1!4,rounded corners=4pt,minimum width=2.8cm,minimum height=0.6cm,font=\tiny\bfseries,text=#1,align=center},
  link/.style={-{Stealth[length=4pt]},#1!50,thick}
]
  \node[level=indigo] (ann) at (0,0) {\faCalendar\ Année Scolaire};
  \node[level=teal, below left=1.2cm and 1cm of ann] (s1) {\faServer\ Semestre 1};
  \node[level=teal, below right=1.2cm and 1cm of ann] (s2) {\faServer\ Semestre 2};
  \node[child=violet, below left=0.8cm and 0.2cm of s1] (p1) {\faCalendar\ Période 1};
  \node[child=violet, below right=0.8cm and 0.2cm of s1] (p2) {\faCalendar\ Période 2};
  \node[child=violet, below left=0.8cm and 0.2cm of s2] (p3) {\faCalendar\ Période 3};
  \node[child=violet, below right=0.8cm and 0.2cm of s2] (p4) {\faCalendar\ Période 4};
  \node[child=amber, below=0.6cm of p1] (e1) {\faPen\ Évals P1};
  \node[child=amber, below=0.6cm of p2] (e2) {\faPen\ Évals P2};
  \node[child=amber, below=0.6cm of p3] (e3) {\faPen\ Évals P3};
  \node[child=amber, below=0.6cm of p4] (e4) {\faPen\ Évals P4};
  \draw[link=indigo] (ann) -- (s1);
  \draw[link=indigo] (ann) -- (s2);
  \draw[link=teal] (s1) -- (p1);
  \draw[link=teal] (s1) -- (p2);
  \draw[link=teal] (s2) -- (p3);
  \draw[link=teal] (s2) -- (p4);
  \draw[link=violet] (p1) -- (e1);
  \draw[link=violet] (p2) -- (e2);
  \draw[link=violet] (p3) -- (e3);
  \draw[link=violet] (p4) -- (e4);
\end{tikzpicture}
\end{center}

\newpage
% ===================================================================
\chapter{Onglet 1 — Types d'Évaluations}
% ===================================================================

\section{Objectif de cet onglet}

L'onglet \onglet{Types d'Évaluations} est le point de départ de tout le processus d'évaluation. Il permet de définir le \textbf{catalogue} des catégories d'évaluations que votre établissement utilisera tout au long de l'année scolaire.

\begin{info}
Les types d'évaluations sont \textbf{partagés au niveau national} via le Hub. Cela signifie que tous les établissements d'un même pays disposent du même catalogue de base. Cependant, chaque établissement peut créer ses propres types additionnels.
\end{info}

\subsection{Présentation de l'interface}

Lorsque vous cliquez sur l'onglet \textbf{Types d'Évaluations}, vous voyez~:

\begin{enumerate}[leftmargin=1.2cm]
  \item \textbf{Un bouton \bouton{+ Nouveau Type}} en haut à droite pour ajouter un nouveau type.
  \item \textbf{Une grille de cartes} affichant tous les types existants. Chaque carte contient~:
  \begin{itemize}
    \item Le \textbf{nom complet} du type (ex.\ : \textit{Interrogation Écrite}).
    \item Le \textbf{sigle} en badge coloré (ex.\ : \textit{IE}).
    \item La \textbf{description} si renseignée.
    \item Deux boutons d'action~: \textcolor{indigo}{\faPen} (modifier) et \textcolor{rose}{\faTrash} (supprimer).
  \end{itemize}
\end{enumerate}

\begin{center}
\begin{tikzpicture}
  % Card frame
  \fill[white, rounded corners=8pt, drop shadow={shadow xshift=1pt, shadow yshift=-1pt, opacity=0.1}] (0,0) rectangle (7,2.5);
  \draw[indigo!20, rounded corners=8pt, line width=1pt] (0,0) rectangle (7,2.5);
  % Top accent bar
  \fill[indigo, rounded corners=8pt] (0,2.2) rectangle (7,2.5);
  % Content
  \node[anchor=west, font=\small\bfseries\color{slate}] at (0.4,1.7) {Interrogation Écrite};
  \node[fill=indigo!12, rounded corners=3pt, inner sep=3pt, font=\tiny\bfseries\color{indigo}] at (5.5,1.7) {IE};
  \node[anchor=west, font=\tiny\color{slatelight}, text width=5cm] at (0.4,1.1) {Interrogation surprise ou planifiée pour évaluer la compréhension du cours.};
  % Action buttons
  \node[fill=indigo!8, rounded corners=3pt, inner sep=3pt, font=\tiny\color{indigo}] at (5.2,0.4) {\faPen};
  \node[fill=rose!8, rounded corners=3pt, inner sep=3pt, font=\tiny\color{rose}] at (6.2,0.4) {\faTrash};
\end{tikzpicture}
\captionof{figure}{Exemple d'une carte de type d'évaluation}
\end{center}

\section{Pour créer un nouveau type d'évaluation}

\begin{prerequis}
Vous devez être connecté avec un compte ayant les droits de gestion des évaluations (directeur ou préfet d'études).
\end{prerequis}

\begin{etape}[title=Procédure détaillée : Créer un type d'évaluation]
\begin{enumerate}[leftmargin=1.5cm,label=\textbf{\arabic*.},itemsep=12pt]
  \item \textbf{Accédez à l'onglet} — Cliquez sur l'onglet \onglet{Types d'Évaluations} situé tout à gauche dans la barre des pills. L'onglet actif se distingue par un fond indigo avec un texte blanc.

  \item \textbf{Ouvrez le formulaire} — Repérez le bouton \bouton{+ Nouveau Type} en haut à droite de la grille de cartes. Cliquez dessus. Une fenêtre modale (pop-up) apparaît au centre de l'écran avec un titre \textit{<<~Nouveau type d'évaluation~>>}.

  \item \textbf{Remplissez le champ Nom} — Dans le champ \champ{Nom}, tapez le nom complet du type d'évaluation. Par exemple~: \textit{Travail Pratique}. Ce champ est \textbf{obligatoire} (marqué d'un astérisque *).

  \item \textbf{Remplissez le champ Sigle} — Dans le champ \champ{Sigle}, tapez l'abréviation du type en 2 à 4 lettres. Par exemple~: \textit{TP}. Le système convertit automatiquement votre saisie en \textbf{majuscules}. Ce champ est \textbf{obligatoire}.

  \item \textbf{Remplissez la Description} — Dans le champ \champ{Description}, vous pouvez (optionnellement) ajouter un texte explicatif. Par exemple~: \textit{<<~Évaluation pratique en laboratoire ou sur le terrain.~>>}

  \item \textbf{Enregistrez} — Cliquez sur le bouton \bouton{Enregistrer} en bas de la modale. Le système envoie les données au serveur via l'API et~:
  \begin{itemize}
    \item Un \textbf{toast vert} \textit{<<~Type enregistré !~>>} confirme la réussite.
    \item La modale se ferme automatiquement.
    \item La grille de cartes se recharge et affiche le nouveau type.
  \end{itemize}
\end{enumerate}
\end{etape}

\begin{resultat}
Après cette opération, le nouveau type apparaît sous forme de carte dans la grille. Il sera disponible dans le menu déroulant \champ{Type} lors de la création d'évaluations (onglet 2).
\end{resultat}

\section{Pour modifier un type existant}

\begin{etape}[title=Procédure détaillée : Modifier un type d'évaluation]
\begin{enumerate}[leftmargin=1.5cm,label=\textbf{\arabic*.},itemsep=10pt]
  \item \textbf{Repérez la carte} du type que vous souhaitez modifier dans la grille.

  \item \textbf{Cliquez sur l'icône} \textcolor{indigo}{\faPen} (stylo bleu) située dans la partie basse de la carte, à côté de l'icône corbeille.

  \item \textbf{La modale s'ouvre} avec les champs pré-remplis. Modifiez le \champ{Nom}, le \champ{Sigle} ou la \champ{Description} selon vos besoins.

  \item \textbf{Cliquez sur} \bouton{Enregistrer}. Le type est mis à jour immédiatement et la grille se rafraîchit.
\end{enumerate}
\end{etape}

\section{Pour supprimer un type d'évaluation}

\begin{danger}
La suppression est \textbf{définitive et irréversible}. Si des évaluations sont rattachées à ce type, elles perdront leur catégorisation. Vérifiez toujours avant de supprimer qu'aucune évaluation active n'utilise ce type.
\end{danger}

\begin{etape}[title=Procédure détaillée : Supprimer un type]
\begin{enumerate}[leftmargin=1.5cm,label=\textbf{\arabic*.},itemsep=10pt]
  \item \textbf{Repérez la carte} du type à supprimer.
  \item \textbf{Cliquez sur l'icône} \textcolor{rose}{\faTrash} (corbeille rouge).
  \item \textbf{Une boîte de confirmation} apparaît~: \textit{<<~Êtes-vous sûr de vouloir supprimer ce type d'évaluation ?~>>}
  \item \textbf{Cliquez sur OK} pour confirmer. Le type disparaît de la grille et un toast de confirmation s'affiche.
\end{enumerate}
\end{etape}

\section{Exemples de types couramment utilisés}

\begin{center}
\begin{tabularx}{\textwidth}{llX}
\toprule
\textbf{Sigle} & \textbf{Nom complet} & \textbf{Description recommandée} \\
\midrule
IE & Interrogation Écrite & Interrogation de courte durée (15--30 min) \\
DS & Devoir Surveillé & Évaluation formelle en classe (1--2h) \\
TP & Travail Pratique & Évaluation en laboratoire ou sur terrain \\
DD & Devoir à Domicile & Travail à rendre réalisé à la maison \\
EX & Examen & Épreuve de fin de période ou semestre \\
QCM & Questionnaire à Choix & Test rapide à choix multiples \\
EO & Évaluation Orale & Présentation ou interrogation orale \\
\bottomrule
\end{tabularx}
\end{center}

\newpage
% ===================================================================
\chapter{Onglet 2 — Planifier les Évaluations}
% ===================================================================

\section{Objectif de cet onglet}

L'onglet \onglet{Planifier les Évaluations} est l'endroit où vous créez les \textbf{évaluations concrètes}~: chaque interrogation, chaque devoir, chaque examen. Vous y définissez pour chaque évaluation~: le cours concerné, la période de rattachement, la date, le barème maximal, et éventuellement le document sujet en PDF.

\begin{info}
Cet onglet organise les évaluations dans un \textbf{système d'accordéons à deux niveaux}~:
\begin{itemize}
  \item \textbf{Niveau 1} — Regroupement par \textbf{cours} (en-tête vert avec icône \faBook).
  \item \textbf{Niveau 2} — Regroupement par \textbf{période} au sein de chaque cours (en-tête violet avec icône \faCalendar).
\end{itemize}
Cliquez sur un en-tête d'accordéon pour le déplier/replier.
\end{info}

\subsection{Présentation de l'interface}

L'interface se compose de~:

\begin{enumerate}[leftmargin=1.2cm]
  \item \textbf{Barre de filtres} en haut — avec deux menus déroulants~:
  \begin{itemize}
    \item \champ{Classe} — sélectionnez la classe cible. \textbf{Obligatoire} pour afficher quoi que ce soit.
    \item \champ{Cours} — filtrez par cours spécifique (optionnel ; si vide, tous les cours sont affichés).
  \end{itemize}
  \item \textbf{Bouton} \bouton{+ Nouvelle Évaluation} — pour créer une évaluation.
  \item \textbf{Zone d'accordéons} — affiche les évaluations regroupées.
\end{enumerate}

\begin{center}
\begin{tikzpicture}[
  acc/.style={rounded corners=6pt, minimum width=13cm, minimum height=0.9cm, font=\small, align=left, drop shadow={shadow xshift=0.5pt, shadow yshift=-0.5pt, opacity=0.08}},
  sub/.style={rounded corners=4pt, minimum width=12cm, minimum height=0.7cm, font=\scriptsize, align=left},
  row/.style={minimum width=11.5cm, minimum height=0.55cm, font=\tiny, align=left}
]
  % Cours accordion
  \node[acc, fill=emerald!85, text=white] (a1) at (0,0) {\quad\faBook\ \textbf{Mathématiques} \hfill \scriptsize 5 évaluations \quad \faChevronDown\quad};
  % Periode
  \node[sub, fill=violet!10, text=violet, below=2pt of a1] (s1) {\quad\quad\faCalendar\ \textbf{1\textsuperscript{ère} Période} — Taux 25\% \hfill \scriptsize 3 évals \quad \faChevronDown\quad};
  % Table header
  \node[row, fill=slate!8, text=slate, below=2pt of s1] (h1) {\quad\quad\quad\footnotesize\textbf{Titre} \hfill \textbf{Type} \hfill \textbf{Date} \hfill \textbf{Max} \hfill \textbf{Actions}\quad};
  % Rows
  \node[row, fill=white, text=slate, below=1pt of h1] (r1) {\quad\quad\quad\footnotesize Interro Chapitre 1 \hfill IE \hfill 15/09/2025 \hfill /10 \hfill \textcolor{indigo}{\faPen}\ \textcolor{rose}{\faTrash}\quad};
  \node[row, fill=lightbg, text=slate, below=1pt of r1] (r2) {\quad\quad\quad\footnotesize Devoir Fonctions \hfill DS \hfill 28/09/2025 \hfill /20 \hfill \textcolor{indigo}{\faPen}\ \textcolor{rose}{\faTrash}\quad};
  \node[row, fill=white, text=slate, below=1pt of r2] (r3) {\quad\quad\quad\footnotesize QCM Algèbre \hfill QCM \hfill 05/10/2025 \hfill /10 \hfill \textcolor{indigo}{\faPen}\ \textcolor{rose}{\faTrash}\quad};
  % Second periode
  \node[sub, fill=violet!10, text=violet, below=4pt of r3] (s2) {\quad\quad\faCalendar\ \textbf{2\textsuperscript{ème} Période} — Taux 25\% \hfill \scriptsize 2 évals \quad \faChevronRight\quad};

  % Border
  \draw[emerald!30, rounded corners=6pt, line width=1pt] ([xshift=-0.3cm,yshift=0.2cm]a1.north west) rectangle ([xshift=0.3cm,yshift=-0.2cm]s2.south east);
\end{tikzpicture}
\end{center}

\section{Pour créer une nouvelle évaluation}

\begin{prerequis}
\begin{itemize}
  \item Au moins un \textbf{type d'évaluation} doit exister (onglet 1).
  \item La \textbf{configuration annuelle} (cycles, classes, cours) doit être en place.
  \item Les \textbf{répartitions} (calendrier scolaire) doivent être définies avec des périodes.
\end{itemize}
\end{prerequis}

\begin{etape}[title=Procédure détaillée : Créer une évaluation]
\begin{enumerate}[leftmargin=1.5cm,label=\textbf{\arabic*.},itemsep=12pt]
  \item \textbf{Sélectionnez la classe} — Ouvrez le menu déroulant \champ{Classe} dans la barre de filtres. Choisissez la classe concernée (ex.\ : \textit{3ème Année Secondaire (Sciences)}). Les cours de cette classe se chargent automatiquement.

  \item \textbf{Sélectionnez le cours} (recommandé) — Bien que facultatif, il est conseillé de filtrer par cours dans le menu \champ{Cours} pour un affichage ciblé.

  \item \textbf{Cliquez sur} \bouton{+ Nouvelle Évaluation} — La fenêtre modale s'ouvre avec le formulaire complet.

  \item \textbf{Remplissez le champ Titre} — Donnez un nom clair et descriptif à l'évaluation. Exemple~: \textit{<<~Interrogation — Les fonctions polynômes~>>}.

  \item \textbf{Sélectionnez le Type} — Choisissez dans le menu déroulant \champ{Type} parmi les types créés à l'onglet 1 (IE, DS, TP, etc.).

  \item \textbf{Sélectionnez le Cours} — Si le cours n'est pas déjà sélectionné dans le filtre, choisissez-le dans le menu déroulant \champ{Cours} de la modale.

  \item \textbf{Sélectionnez la Période} — Le menu déroulant \champ{Période} liste toutes les périodes ouvertes de la classe. Choisissez celle dans laquelle cette évaluation a lieu (ex.\ : \textit{1\textsuperscript{ère} Période}).

  \item \textbf{Saisissez la Date de passation} — Tapez la date au format \texttt{jj/mm/aaaa} dans le champ \champ{Date}. Le formatage automatique ajoute les barres obliques au fur et à mesure~: il suffit de taper les 8 chiffres.

  \item \textbf{Saisissez la Date de remise} (optionnel) — Si l'évaluation a une date de remise différente.

  \item \textbf{Définissez le Maximum} — Dans le champ \champ{Maximum (pts)}, entrez le barème maximal (ex.\ : 10, 20, 30...). Ce maximum sera utilisé lors de la saisie des notes.

  \item \textbf{Ajoutez les Consignes} (optionnel) — Rédigez des instructions dans le champ texte \champ{Consignes / Description}.

  \item \textbf{Joignez un document PDF} (optionnel) — Cliquez sur \champ{Document PDF} pour joindre le sujet de l'évaluation.

  \item \textbf{Enregistrez} — Cliquez sur \bouton{Enregistrer}. Un toast vert confirme la réussite.
\end{enumerate}
\end{etape}

\begin{resultat}
L'évaluation apparaît dans l'accordéon du cours concerné, sous la période sélectionnée. Elle est désormais disponible pour la saisie de notes (onglet 3).
\end{resultat}

\section{Pour modifier une évaluation existante}

\begin{etape}[title=Procédure : Modifier une évaluation]
\begin{enumerate}[leftmargin=1.5cm,label=\textbf{\arabic*.},itemsep=8pt]
  \item Dépliez l'accordéon du cours, puis celui de la période.
  \item Repérez l'évaluation dans la liste tabulaire.
  \item Cliquez sur \textcolor{indigo}{\faPen} dans la colonne \textbf{Actions}.
  \item La modale s'ouvre pré-remplie. Modifiez les champs nécessaires.
  \item Cliquez sur \bouton{Enregistrer}.
\end{enumerate}
\end{etape}

\section{Pour supprimer une évaluation}

\begin{attention}
La suppression d'une évaluation supprime aussi \textbf{toutes les notes} qui y sont rattachées. Procédez avec précaution.
\end{attention}

Cliquez sur \textcolor{rose}{\faTrash} dans la colonne Actions et confirmez.

% ===================================================================
\chapter{Onglet 3 --- Notes Par \'Evaluation}
% ===================================================================

\section{Objectif de cet onglet}

L'onglet \onglet{Notes Par \'Evaluation} est l'interface principale de \textbf{saisie des notes brutes}. C'est ici que chaque enseignant ou administrateur entre les r\'esultats obtenus par les \'el\`eves pour chaque \'evaluation planifi\'ee \`a l'onglet~2.

Cet onglet supporte~:
\begin{itemize}
  \item \textbf{Trois modes de saisie} --- Par P\'eriode, Par Cours, Par \'Evaluation.
  \item \textbf{Deux types d'entr\'ee} --- Notes de p\'eriode (\'evaluations continues) et Notes d'examen.
  \item \textbf{Import/Export Excel} --- Pour traiter les notes en lot.
  \item \textbf{Navigation au clavier} --- Pour une saisie rapide et ergonomique.
\end{itemize}

\section{Comprendre l'interface}

L'interface se compose de trois zones principales~:

\begin{enumerate}[leftmargin=1.2cm]
  \item \textbf{Zone de filtres} (en haut) --- Contient les menus d\'eroulants \champ{Classe} et \champ{P\'eriode}. La s\'election d'une classe charge automatiquement les p\'eriodes disponibles.
  \item \textbf{Boutons de mode} --- Trois boutons permettent de choisir le niveau de d\'etail de la grille~: \textbf{Par P\'eriode}, \textbf{Par Cours} ou \textbf{Par \'Evaluation}.
  \item \textbf{Grille de saisie} --- Un tableau interactif avec les \'el\`eves en lignes et les \'evaluations en colonnes. Chaque cellule est un champ de saisie num\'erique.
\end{enumerate}

\section{Pour s\'electionner la classe et la p\'eriode}

\begin{etape}[title=Proc\'edure : Pr\'eparer la grille de saisie]
\begin{enumerate}[leftmargin=1.5cm,label=\textbf{\arabic*.},itemsep=12pt]
  \item \textbf{S\'electionnez une Classe} --- Ouvrez le menu d\'eroulant \champ{Classe} et choisissez la classe. Le syst\`eme charge les r\'epartitions (p\'eriodes) disponibles.

  \item \textbf{S\'electionnez une P\'eriode} --- Le menu d\'eroulant \champ{P\'eriode} affiche la liste organis\'ee en groupes~:
  \begin{itemize}
    \item Les \textbf{p\'eriodes normales} (1\textsuperscript{\`ere} P\'eriode, 2\textsuperscript{\`eme} P\'eriode...) pour les notes continues.
    \item Les entr\'ees \textbf{\faPencil\ Examen} (en violet dans le menu) pour les notes d'examen semestrielles.
  \end{itemize}

  \item \textbf{La grille se charge} automatiquement avec les \'el\`eves de la classe et les \'evaluations de la p\'eriode.
\end{enumerate}
\end{etape}

\begin{attention}
Si vous s\'electionnez une entr\'ee \textbf{Examen}, la grille passe en \textbf{mode examen}~: les colonnes affichent les cours (pas les \'evaluations individuelles), et le maximum correspond au bar\`eme d'examen configur\'e pour chaque cours.
\end{attention}

\section{Les trois modes de saisie en d\'etail}

Apr\`es le chargement de la grille, trois boutons de mode apparaissent au-dessus du tableau~:

\subsection{Mode \guillemotleft~Par P\'eriode~\guillemotright}

\begin{etape}[title=Mode Par P\'eriode --- Vue panoramique]
C'est le mode \textbf{par d\'efaut}. Il affiche un tableau large avec~:
\begin{itemize}
  \item \textbf{Lignes} --- tous les \'el\`eves de la classe, tri\'es alphab\'etiquement.
  \item \textbf{Colonnes} --- \textit{toutes} les \'evaluations de \textit{tous} les cours pour la p\'eriode s\'electionn\'ee.
  \item \textbf{En-t\^etes} --- regroup\'es par cours (en-t\^ete de groupe color\'e) puis par \'evaluation.
\end{itemize}

\textbf{Utilisation recommand\'ee}~: Quand vous voulez avoir une vue d'ensemble sur toutes les notes d'une p\'eriode.
\end{etape}

\subsection{Mode \guillemotleft~Par Cours~\guillemotright}

\begin{etape}[title=Mode Par Cours --- Vue cibl\'ee]
Ce mode ajoute un filtre \champ{Cours} suppl\'ementaire. Il affiche uniquement les \'evaluations d'\textbf{un seul cours}.

\textbf{Pour activer}~: cliquez sur le bouton \btnviolet{Par Cours}, puis s\'electionnez le cours dans le menu d\'eroulant qui appara\^it.

\textbf{Utilisation recommand\'ee}~: Quand un enseignant veut saisir les notes de son cours uniquement.
\end{etape}

\subsection{Mode \guillemotleft~Par \'Evaluation~\guillemotright}

\begin{etape}[title=Mode Par \'Evaluation --- Vue unitaire]
Ce mode est le plus cibl\'e. Il affiche une \textbf{seule colonne} pour une seule \'evaluation d'un seul cours.

\textbf{Pour activer}~: cliquez sur \btnviolet{Par \'Evaluation}, s\'electionnez le cours, puis l'\'evaluation sp\'ecifique.

\textbf{Utilisation recommand\'ee}~: Pour une saisie rapide et concentr\'ee apr\`es une interrogation.
\end{etape}

\section{Pour saisir les notes}

\begin{etape}[title=Proc\'edure d\'etaill\'ee : Saisir les notes]
\begin{enumerate}[leftmargin=1.5cm,label=\textbf{\arabic*.},itemsep=12pt]
  \item \textbf{Cliquez sur une cellule} de la grille. Le curseur se place dans le champ de saisie. Le fond de la cellule devient l\'eg\`erement bleut\'e pour indiquer la s\'election.

  \item \textbf{Tapez la note} sous forme num\'erique (entier ou d\'ecimal avec un point). Par exemple~: \texttt{15} ou \texttt{7.5}.

  \item \textbf{Observez la validation visuelle}~:
  \begin{itemize}
    \item \textcolor{emerald}{\faCheck\ Fond vert clair} --- la note est dans les limites (entre 0 et le maximum).
    \item \textcolor{rose}{\faTimes\ Fond rouge clair} --- la note d\'epasse le maximum autoris\'e. Elle sera refus\'ee \`a l'enregistrement.
  \end{itemize}

  \item \textbf{Naviguez avec le clavier}~:
  \begin{itemize}
    \item \texttt{Tab} ou \texttt{Entr\'ee} --- passe \`a la cellule \textbf{suivante} (vers la droite, puis ligne suivante).
    \item \texttt{Shift+Tab} --- revient \`a la cellule \textbf{pr\'ec\'edente}.
    \item \texttt{Fl\`eche Bas} --- descend dans la \textbf{m\^eme colonne}.
    \item \texttt{Fl\`eche Haut} --- monte dans la m\^eme colonne.
  \end{itemize}

  \item \textbf{Enregistrez} --- Une fois toutes les notes saisies, cliquez sur \bouton{Enregistrer} (en bas ou en haut de la grille). Le syst\`eme envoie toutes les notes au serveur et affiche un toast de confirmation avec le nombre de notes enregistr\'ees.
\end{enumerate}
\end{etape}

\begin{astuce}
Vous n'\^etes pas oblig\'e de remplir toutes les cellules avant d'enregistrer. Les cellules vides seront simplement ignor\'ees. Vous pourrez compl\'eter plus tard.
\end{astuce}

\section{Pour saisir les notes d'examen}

\begin{etape}[title=Proc\'edure : Saisir les notes d'examen]
\begin{enumerate}[leftmargin=1.5cm,label=\textbf{\arabic*.},itemsep=10pt]
  \item \textbf{S\'electionnez la classe} dans le filtre.
  \item \textbf{Dans le menu P\'eriode}, choisissez l'entr\'ee \textbf{\faPencil\ Examen --- [Nom du semestre]} (affich\'ee en violet avec une ic\^one stylo).
  \item La grille d'examen se charge~: les colonnes sont les \textbf{cours} (pas les \'evaluations individuelles).
  \item Saisissez les notes d'examen cours par cours pour chaque \'el\`eve.
  \item Cliquez sur \bouton{Enregistrer Examen}.
\end{enumerate}
\end{etape}

\section{Pour importer des notes depuis Excel}

\begin{etape}[title=Proc\'edure : Import Excel]
\begin{enumerate}[leftmargin=1.5cm,label=\textbf{\arabic*.},itemsep=10pt]
  \item \textbf{T\'el\'echargez le mod\`ele} --- Cliquez sur \btnamber{Mod\`ele Excel}. Un fichier \texttt{.xlsx} se t\'el\'echarge avec la liste des \'el\`eves et les en-t\^etes des \'evaluations.
  \item \textbf{Remplissez le fichier} --- Ouvrez-le dans Excel ou LibreOffice Calc, saisissez les notes dans les colonnes correspondantes.
  \item \textbf{Importez} --- Cliquez sur \btnamber{Importer}, s\'electionnez votre fichier.
  \item Les notes sont charg\'ees directement dans la grille. V\'erifiez et cliquez sur \bouton{Enregistrer}.
\end{enumerate}
\end{etape}

\newpage
% ===================================================================
\chapter{Onglet 4 --- Notes Pond\'er\'ees}
% ===================================================================

\section{Objectif de cet onglet}

L'onglet \onglet{Notes Pond\'er\'ees} est le \textbf{centre de calcul} du bulletin scolaire. Il transforme les notes brutes saisies \`a l'onglet~3 en r\'esultats agr\'eg\'es~: \textbf{NP} (Note Pond\'er\'ee par p\'eriode), \textbf{TJ} (Travaux Journaliers) et \textbf{TOTAL} (TJ + EX).

\subsection{Pipeline de calcul}

Le calcul suit un pipeline en 5 \'etapes~:

\begin{center}
\begin{tikzpicture}[
  box/.style={draw=slate!40, fill=white, rounded corners=6pt, minimum width=2cm, minimum height=1cm, align=center, font=\scriptsize\bfseries, text=slate, drop shadow={shadow xshift=0.5pt, shadow yshift=-0.5pt, opacity=0.1}},
  arrow/.style={-{Stealth[length=5pt,width=4pt]}, thick, color=slatelight},
  lbl/.style={font=\tiny\itshape\color{slatelight}, above, midway}
]
  \node[box] (e1) at (0,0) {\faPencil\\[2pt]Notes\\brutes};
  \node[box, right=0.4cm of e1] (np) {\faCalculator\\[2pt]Note\\Pond\'er\'ee};
  \node[box, right=0.4cm of np] (tj) {\faServer\\[2pt]TJ\\Semestriel};
  \node[box, right=0.4cm of tj] (ex) {\faFile\\[2pt]EX\\Examen};
  \node[box, right=0.4cm of ex] (tot) {\faTrophy\\[2pt]TOTAL};
  \draw[arrow] (e1) -- (np) node[lbl]{poids};
  \draw[arrow] (np) -- (tj) node[lbl]{somme};
  \draw[arrow] (tj) -- (tot);
  \draw[arrow] (ex) -- (tot) node[lbl]{TJ+EX};
\end{tikzpicture}
\end{center}

\section{Comprendre la vue hi\'erarchique}

Lorsque vous s\'electionnez une classe, l'interface affiche un \textbf{accord\'eon imbriqu\'e \`a 3 niveaux}~:

\begin{etape}[title=Structure de l'accord\'eon]
\begin{description}[leftmargin=1.5cm, style=nextline, font=\bfseries\color{indigo}]
  \item[Niveau 1 --- Cours] Chaque cours est un bloc color\'e avec le nom du cours et un bouton \textcolor{amber}{\faBolt\ TJ} pour calculer les TJ en lot. Cliquez pour d\'eplier et voir les semestres.
  \item[Niveau 2 --- Semestre/Trimestre] \`A l'int\'erieur de chaque cours, les semestres sont affich\'es avec fond pastel. Ils montrent le nombre de p\'eriodes enfants. Cliquez pour d\'eplier le tableau.
  \item[Niveau 3 --- Tableau de donn\'ees] Le tableau affiche pour chaque \'el\`eve~:
  \begin{itemize}
    \item Les \textbf{notes brutes} de chaque \'evaluation (fond jaune clair).
    \item La \textbf{NP} (Note Pond\'er\'ee) par p\'eriode (fond orange, bordure droite orange).
    \item Le \textbf{TJ} semestriel (fond vert clair).
    \item L'\textbf{EX} examen (fond violet clair).
    \item Le \textbf{TOTAL} = TJ + EX (fond bleu clair, en gras).
  \end{itemize}
\end{description}
\end{etape}

\section{Pour calculer les TJ d'un cours (pond\'eration par p\'eriode)}

\begin{etape}[title=Proc\'edure d\'etaill\'ee : Calculer TJ avec pond\'eration]
\begin{enumerate}[leftmargin=1.5cm,label=\textbf{\arabic*.},itemsep=12pt]
  \item \textbf{Rep\'erez le cours} dans l'accord\'eon et cliquez sur le bouton \textcolor{amber}{\faBolt\ TJ} dans l'en-t\^ete du cours.

  \item \textbf{La modale de pond\'eration s'ouvre}~: elle affiche pour chaque semestre un bloc color\'e contenant ses p\'eriodes, et pour chaque p\'eriode, la liste de ses \'evaluations sous forme de tableau.

  \item \textbf{Pour chaque \'evaluation}, vous voyez~:
  \begin{itemize}
    \item Une \textbf{case \`a cocher} --- pour inclure/exclure l'\'evaluation du calcul.
    \item Le \textbf{nom} de l'\'evaluation et son \textbf{maximum} (/10, /20...).
    \item Le \textbf{poids en \%} --- modifiable. Par d\'efaut, le poids est r\'eparti \'equitablement entre toutes les \'evaluations coch\'ees.
  \end{itemize}

  \item \textbf{Ajustez les poids} --- Vous pouvez modifier manuellement le pourcentage de chaque \'evaluation. Si vous d\'ecochez une \'evaluation, les poids se \textbf{redistribuent automatiquement} entre les \'evaluations restantes (redistribution \'equitable).

  \item \textbf{Lancez le calcul} \`a l'un des trois niveaux~:
  \begin{itemize}
    \item \textbf{Calc.} (petit bouton par p\'eriode) --- calcule la NP d'une seule p\'eriode.
    \item \textbf{Calculer} (bouton par semestre) --- calcule toutes les p\'eriodes du semestre.
    \item \textbf{Calculer TOUT} (bouton global en bas) --- calcule toutes les p\'eriodes de tous les semestres pour ce cours.
  \end{itemize}

  \item \textbf{Le tableau se met \`a jour} avec les valeurs NP calcul\'ees pour chaque \'el\`eve.
\end{enumerate}
\end{etape}

\begin{astuce}
Lorsque vous cochez/d\'ecochez une \'evaluation, les poids se redistribuent automatiquement. Par exemple, avec 4 \'evaluations coch\'ees, chacune re\c{c}oit 25\%. Si vous en d\'ecochez une, les 3 restantes passent \`a 33,33\% chacune.
\end{astuce}

\section{Pour calculer le R\'esultat semestriel (TJ + EX $\rightarrow$ TOTAL)}

\begin{etape}[title=Proc\'edure : Calculer TJ + EX = TOTAL]
\begin{enumerate}[leftmargin=1.5cm,label=\textbf{\arabic*.},itemsep=10pt]
  \item \textbf{D\'epliez} l'accord\'eon d'un cours puis d'un semestre pour voir le tableau.
  \item \textbf{Rep\'erez} le groupe de colonnes \textbf{R\'esultat} (TJ, EX, TOTAL) \`a droite du tableau.
  \item \textbf{Cliquez} sur l'ic\^one \textcolor{amber}{\faCalculator} dans l'en-t\^ete du groupe R\'esultat.
  \item \textbf{La modale de confirmation} appara\^it et explique le processus en 3 \'etapes~:
  \begin{enumerate}
    \item \textbf{TJ} --- Le syst\`eme somme les NP de toutes les p\'eriodes enfants, pond\'er\'ees par leur taux de participation.
    \item \textbf{EX} --- Il agr\`ege les notes d'examen saisies au niveau parent.
    \item \textbf{TOTAL} --- Il additionne TJ + EX.
  \end{enumerate}
  \item \textbf{Cliquez sur} \bouton{Calculer}. Les colonnes TJ, EX et TOTAL se remplissent.
\end{enumerate}
\end{etape}

\section{Pour calculer tout en lot (tous les cours)}

Le bouton \bouton{Calculer Tout} situ\'e en haut de la page lance le calcul pour \textbf{tous les cours} de la classe en une seule op\'eration. Le syst\`eme it\`ere sur chaque semestre/trimestre parent et calcule les TJ + EX $\rightarrow$ TOTAL pour chaque cours.

\begin{resultat}
Apr\`es le calcul global, un compteur indique le nombre total de notes calcul\'ees (ex.~: \textit{\guillemotleft~245 note(s) calcul\'ee(s)~\guillemotright}). Le tableau se rafra\^ichit automatiquement.
\end{resultat}

% ===================================================================
\chapter{Sections compl\'ementaires}
% ===================================================================

\section{D\'elib\'erations}

La section \textbf{D\'elib\'erations} (accessible via l'ic\^one \faGavel) permet de valider officiellement les r\'esultats scolaires. Elle supporte 4 niveaux de d\'elib\'eration~: \textbf{Par P\'eriode}, \textbf{Par Trimestre/Semestre}, \textbf{Par Ann\'ee} et \textbf{Rep\^echage}.

\begin{etape}[title=Proc\'edure : Lancer une d\'elib\'eration]
\begin{enumerate}[leftmargin=1.5cm,label=\textbf{\arabic*.},itemsep=10pt]
  \item \textbf{Choisissez le type} en cliquant sur l'onglet correspondant (P\'eriode, Trimestre, Ann\'ee, Rep\^echage).
  \item \textbf{S\'electionnez une Classe} dans le menu d\'eroulant.
  \item \textbf{S\'electionnez la R\'epartition} (pour P\'eriode/Trimestre) ou la \textbf{Session} (pour Ann\'ee/Rep\^echage).
  \item Un panneau de \textbf{Conditions de D\'elib\'eration} s'affiche montrant les seuils, mentions et sanctions configur\'es.
  \item \textbf{Cliquez sur} \bouton{Lancer la D\'elib\'eration}.
  \item Le badge de statut passe \`a \textit{\guillemotleft~D\'elib\'eration en cours...\guillemotright} puis \`a \textit{\guillemotleft~R\'eussie~\guillemotright}.
  \item Un \textbf{tableau de r\'esultats} appara\^it avec les colonnes~:
  \begin{itemize}
    \item \textbf{\#} --- classement (m\'edailles d'or/argent/bronze pour les 3 premiers).
    \item \textbf{\'El\`eve} --- nom et pr\'enom.
    \item \textbf{Total} --- note totale / maximum.
    \item \textbf{\%} --- pourcentage (color\'e vert $\geq$70\%, orange $\geq$50\%, rouge $<$50\%).
    \item \textbf{Place} --- rang dans la classe.
    \item \textbf{Mention} --- selon les conditions configur\'ees.
    \item \textbf{D\'ecision} --- Admis (badge vert) ou Ajourn\'e (badge rouge).
  \end{itemize}
\end{enumerate}
\end{etape}

\begin{danger}
\textbf{Annulation en cascade}~: Si vous annulez une d\'elib\'eration de p\'eriode, toutes les d\'elib\'erations de niveaux sup\'erieurs (trimestre, ann\'ee) qui en d\'ependent seront aussi annul\'ees. Le syst\`eme affiche un avertissement d\'etaill\'e avant de proc\'eder.
\end{danger}

\section{Bulletins PDF}

\begin{etape}[title=Proc\'edure : G\'en\'erer des bulletins]
\begin{enumerate}[leftmargin=1.5cm,label=\textbf{\arabic*.},itemsep=10pt]
  \item \textbf{S\'electionnez une classe d\'elib\'er\'ee} dans le menu d\'eroulant de la section Bulletins. Seules les classes ayant au moins une d\'elib\'eration valid\'ee apparaissent.
  \item \textbf{La liste des \'el\`eves} s'affiche sous forme de pastilles cliquables.
  \item \textbf{Cochez les \'el\`eves} souhait\'es en cliquant sur les pastilles (ou cliquez \bouton{Tous}).
  \item Le compteur en bas indique le nombre de s\'electionn\'es (ex.~: \textit{\guillemotleft~\textbf{12} s\'electionn\'e(s)~\guillemotright}).
  \item \textbf{Cliquez sur} \bouton{G\'en\'erer}. Le PDF s'ouvre dans un nouvel onglet du navigateur.
\end{enumerate}
\end{etape}

% ===================================================================
\chapter{R\'ef\'erence rapide}
% ===================================================================

\section{R\'ecapitulatif des endpoints API}

{\small
\begin{tabularx}{\textwidth}{lX}
\toprule
\textbf{Endpoint} & \textbf{Description} \\
\midrule
\texttt{/api/evaluations/types/} & GET --- Liste tous les types d'\'evaluations \\
\texttt{/api/evaluations/types/save/} & POST --- Cr\'eer ou modifier un type \\
\texttt{/api/evaluations/types/delete/} & POST --- Supprimer un type \\
\texttt{/api/evaluations/list/} & GET --- Liste les \'evaluations d'une classe/cours \\
\texttt{/api/evaluations/save/} & POST --- Cr\'eer ou modifier une \'evaluation \\
\texttt{/api/evaluations/delete/} & POST --- Supprimer une \'evaluation \\
\texttt{/api/notes/grid/} & GET --- Grille de notes d'une p\'eriode \\
\texttt{/api/notes/save/} & POST --- Enregistrer les notes saisies \\
\texttt{/api/notes/exam/grid/} & GET --- Grille des notes d'examen \\
\texttt{/api/notes/exam/save/} & POST --- Enregistrer les notes d'examen \\
\texttt{/api/notes/template/} & GET --- Mod\`ele Excel pr\'e-rempli \\
\texttt{/api/notes/import/} & POST --- Importer notes depuis Excel \\
\texttt{/api/notes/bulletin/overview/} & GET --- Vue d'ensemble Notes Pond\'er\'ees \\
\texttt{/api/notes/bulletin/calculate/} & POST --- Calculer TJ + EX = TOTAL \\
\texttt{/api/notes/bulletin/calculate-period-batch/} & POST --- Calcul TJ multi-p\'eriodes \\
\texttt{/api/deliberations/execute/} & POST --- Lancer une d\'elib\'eration \\
\texttt{/api/deliberations/cancel/} & POST --- Annuler une d\'elib\'eration \\
\texttt{/api/deliberations/results/} & GET --- R\'esultats d'une d\'elib\'eration \\
\bottomrule
\end{tabularx}
}

\section{Codes couleur dans les grilles}

\begin{tabularx}{\textwidth}{p{1.5cm}lX}
\toprule
\textbf{Couleur} & \textbf{\'El\'ement} & \textbf{Signification} \\
\midrule
Jaune & \'Evaluations & Notes brutes des interrogations et devoirs \\
Orange & NP & Note Pond\'er\'ee (calcul\'ee par p\'eriode) \\
Vert & TJ & Travaux Journaliers (somme des NP) \\
Violet & EX & Note d'examen \\
Bleu & TOTAL & R\'esultat final (TJ + EX) \\
\bottomrule
\end{tabularx}

\section{Raccourcis clavier}

\begin{tabularx}{\textwidth}{lX}
\toprule
\textbf{Touche} & \textbf{Action} \\
\midrule
\texttt{Tab} / \texttt{Entr\'ee} & Passer \`a la cellule suivante \\
\texttt{Shift + Tab} & Revenir \`a la cellule pr\'ec\'edente \\
\texttt{Fl\`eche Bas} & Descendre dans la m\^eme colonne \\
\texttt{Fl\`eche Haut} & Monter dans la m\^eme colonne \\
\texttt{Echap} & Fermer une fen\^etre modale \\
\bottomrule
\end{tabularx}

\vspace{2cm}
\begin{center}
\begin{tikzpicture}
  \fill[coverbg, rounded corners=12pt] (0,0) rectangle (14,2.5);
  \node[text=white, font=\large\bfseries] at (7,1.5) {\faCheck\hspace{8pt}Fin du Manuel --- Module \'Evaluations};
  \node[text=white!50, font=\small] at (7,0.7) {Plateforme MonEcole --- Nexora Tech --- Avril 2026 --- v1.0};
\end{tikzpicture}
\end{center}


\end{document}
